\documentclass[11pt]{article}

% ─── Packages ────────────────────────────────────────────────────────────────
\usepackage[margin=1in]{geometry}
\usepackage{amsmath,amssymb}
\usepackage{booktabs}
\usepackage{multirow}
\usepackage{graphicx}
\usepackage{xcolor}
\usepackage[colorlinks=true,linkcolor=black,citecolor=black,urlcolor=blue]{hyperref}
\usepackage{natbib}
\usepackage{microtype}
\usepackage{enumitem}
\usepackage{caption}
\usepackage{subcaption}
\usepackage{algorithm}
\usepackage{algorithmic}
\usepackage{url}
\usepackage{array}
\usepackage{tabularx}

% ─── Typography ───────────────────────────────────────────────────────────────
\usepackage[T1]{fontenc}
\usepackage{lmodern}
\renewcommand{\ttdefault}{cmtt}

% ─── Custom commands ─────────────────────────────────────────────────────────
\newcommand{\armi}{\textsc{Arm\,0}}
\newcommand{\armii}{\textsc{Arm\,2}}
\newcommand{\latch}{\textsc{Latch}}
\newcommand{\svss}{\textsc{svss}}
\newcommand{\mcnem}{McNemar}
\newcommand{\pval}[1]{$p\!=\!#1$}
\newcommand{\chisq}[1]{$\chi^2(1)\!=\!#1$}

\definecolor{resultblue}{HTML}{2563EB}
\definecolor{resultamber}{HTML}{D97706}

% ─── Title block ─────────────────────────────────────────────────────────────
\title{When Gates Abstain Too Often: Typed Operator Gates\\
       Fail to Beat Chain-of-Thought on Multi-Hop QA}

\author{Kristian Krattiger\\
        Independent Researcher\\
        \texttt{kristiankratt@gmail.com}}

\date{November 2026}

\begin{document}
\maketitle

% ─── Abstract ─────────────────────────────────────────────────────────────────
\begin{abstract}
We report a negative result from the \latch{} pilot, a preregistered study
comparing a typed-operator gate (Arm~2) against direct chain-of-thought
generation (Arm~0) on the 2WikiMultihopQA decomposable-question subset.
Arm~2 formalises each question as a typed JSON step schema and admits
it only when provenance, schema conformance, and operator-type consistency
all pass.  On $n{=}72$ manually scored questions, \armi{} achieves
80.6\% accuracy at full (100\%) coverage.  \armii{} achieves only 37.5\%
full-coverage accuracy ($-43.1$ pp) because it abstains on 55.6\% of
questions; among the 32 questions it does admit, precision reaches 84.4\%.
A preregistered \mcnem{} test on the jointly-answered subset
yields \chisq{0.67}, \pval{0.41}—far from the \mbox{+15 pp} and $p<0.05$
bar required for promotion.  Per-type analysis shows the abstain
failure is total on \emph{compositional} questions (92\% abstain, 0\%
precision on 2 admits, 0\% full-coverage accuracy), and partial on
\emph{bridge\_comparison} and \emph{comparison} types (68–59\% coverage,
88–92\% precision when admitted).  We identify high abstention, not low
precision, as the decisive failure mode, and discuss implications for
hard symbolic gates in a neurosymbolic multi-hop pipeline.
\end{abstract}

% ─── 1. Introduction ──────────────────────────────────────────────────────────
\section{Introduction}
\label{sec:intro}

Multi-hop question answering requires chaining several inference steps whose
individual correctness can in principle be verified independently.  This
structural property motivates neurosymbolic approaches
\citep{logicLM,LINC,satlm,faithfulCoT,proofOfThought}: rather than relying on
a language model to chain all steps implicitly, one can require the model to
\emph{formalise} its reasoning as a typed intermediate representation,
then \emph{verify} each step mechanically before accepting the answer.

The appeal is clear.  If formalisation succeeds, a symbolic verifier can
catch semantic errors that a soft evaluator might miss.  The two-module
structure---formaliser plus gate---also separates concerns: the gate can be
improved independently of the LLM, and its decisions are auditable.

The \latch{} pilot tests whether this advantage materialises in practice on
2WikiMultihopQA \citep{2wikimultihopqa}.  We define:

\begin{itemize}[leftmargin=*,noitemsep]
  \item \textbf{\armi{} (baseline)}: direct chain-of-thought generation,
        no gate.
  \item \textbf{\armii{} (treatment)}: LLM formaliser (Qwen2.5-14B with
        few-shot examples) produces a typed JSON step schema; a symbolic gate
        checks provenance, schema validity, and operator-type consistency.
        Only questions that pass all checks are admitted; the rest abstain.
\end{itemize}

We preregistered a promotion criterion of $+15$ pp full-coverage accuracy
\emph{and} $p < 0.05$ on a \mcnem{} test.  The gate does not meet either
criterion and the pilot is killed.

This paper makes the following contributions:
\begin{enumerate}[leftmargin=*,noitemsep]
  \item A clean negative result under a preregistered protocol, filling a
        gap in the neurosymbolic literature where positive results are
        over-represented.
  \item An empirical diagnosis of the failure: high abstention on
        compositional question types, not low precision on admitted ones.
  \item A coverage--precision analysis showing the gate's tradeoff point
        (44.4\% coverage, 15.6\% risk) versus the baseline (100\% coverage,
        19.4\% risk).
  \item A discussion of when hard symbolic gates can and cannot be
        expected to help on multi-hop QA.
\end{enumerate}

% ─── 2. Related Work ──────────────────────────────────────────────────────────
\section{Related Work}
\label{sec:related}

\paragraph{Neurosymbolic single-model reasoning.}
Several systems address reasoning by coupling an LLM formaliser with a
symbolic solver.  Logic-LM \citep{logicLM} maps natural language to a
symbolic program, executes it with a solver, and iterates with
self-refinement, gaining $+39.2$\% over plain LLM and $+18.4$\% over CoT
on logical-reasoning benchmarks.  LINC \citep{LINC} translates to first-order
logic and uses Prover9; on ProofWriter, a 15.5B open model plus LINC beats
GPT-3.5/GPT-4 CoT.  SatLM \citep{satlm} emits SAT/SMT specifications
solved by Z3, achieving state-of-the-art on LSAT and BoardgameQA.
Faithful CoT \citep{faithfulCoT} compiles natural language to a symbolic
chain before invoking a solver, outperforming CoT on 9 of 10 benchmarks.
Proof of Thought \citep{proofOfThought} uses a JSON DSL mapped to FOL and
verified by Z3---the closest analogue to \latch{}'s JSON step schema plus
gate.  Programme-of-thought approaches (PAL \citep{pal}, PoT \citep{pot})
let the LLM write Python executed by an interpreter, gaining $+15$ pp on
GSM8K.  DecomP \citep{decomp} decomposes questions into typed sub-tasks
dispatched to a handler library, the closest precedent for typed handler
dispatch.

Notably, \citet{robustnessNeurosym} find that neurosymbolic reasoners are
more robust to perturbation but sometimes score lower overall than pure
LLM baselines---a genuine negative result that partially anticipates ours.

\paragraph{Multi-agent coordination and failure.}
Multi-agent debate (MAD) literature has consistently found it difficult to
beat simple baselines: \citet{goingMAD} show MAD does not reliably outperform
self-consistency; \citet{ifMAD} confirm across five MAD methods and nine
benchmarks.  MAST \citep{mast} analyses 14 multi-agent system failure modes
and finds 79\% arise from specification and coordination problems, not raw
capability---consistent with our finding that the gate's high abstention
rate reflects a coordination mismatch between formaliser and verifier.
ORCH \citep{orch} reports a deterministic orchestration approach that does
beat single models (0.94 vs.\ 0.44 on GSM8K), suggesting structural gains
are possible but highly sensitive to task type.

\paragraph{Gate mechanisms: hard versus soft.}
DebateCoder \citep{debateCoder} applies a soft 95\%-confidence gate to
Pangu-1B code generation, achieving 70.12\% Pass@1 on HumanEval.  VPR
\citep{vpr} uses symbolic oracles as an RL training signal (not an
inference gate).  The survey by \citet{trustButVerify} identifies hard
vs.\ soft gate as a named design axis; \citet{generateFilter} explicitly
names ``hard semantic gates'' in the RLVR literature.  Our work provides
empirical data on how a hard gate degrades coverage on a natural-language
task where the formalisation itself is generated by an LLM.

\paragraph{Error correlation.}
Since \armi{} and \armii{} share the same backbone LLM (the formaliser is
Qwen2.5-14B; the CoT generator is also Qwen-family), correlated errors
are expected.  \citet{llmsJury} measure within-model error correlation at
$\bar\rho = 0.68$, same-family at 0.52.  \citet{correlatedErrors} find
60\% agreement when two LLMs both err.  \citet{nineJudges} show nine
frontier LLMs provide only about two effective independent votes on NLI.
These findings explain the tight \mcnem{} concordance table we observe
(§\ref{sec:results}).

\paragraph{Step-level deductive verification.}
Natural Program \citep{naturalProgram} decomposes CoT into minimal-premise
steps and uses an LM to judge entailment---a soft analogue to our hard
mechanical gate.  \citet{swrv} apply a Z3 verifier to Llama-3-8B CoT,
reporting 0.078 F1 and 72\% solver success---confirming that mechanical
verification of LLM-generated proofs is fragile.  The LLMSR@XLLM25
shared task \citep{llmsrTask,llmsrFifth,llmsrThird} provides further
evidence from the structured reasoning community.

% ─── 3. Gate Design ───────────────────────────────────────────────────────────
\section{The \latch{} Gate Design}
\label{sec:design}

\subsection{Step Schema}

The formaliser maps each question to a \emph{step schema}: a JSON object
listing a sequence of typed operator invocations.  Each step has a type,
input references (either direct evidence spans or prior step outputs), and
an expected output.  The four operators are:

\begin{description}[leftmargin=*,noitemsep]
  \item[\texttt{entity\_equal}] Checks whether two entity references refer
    to the same real-world entity.
  \item[\texttt{temporal\_order}] Verifies a before/after claim between
    two dated events, grounded in evidence spans.
  \item[\texttt{arithmetic}] Performs a numeric calculation over values
    extracted from spans, returning a number or comparison.
  \item[\texttt{conjoin}] Combines two or more atomic results into a
    composite claim, allowing multi-hop chains.
\end{description}

Evidence spans are required to be verbatim substrings of the referenced
document (\emph{text-search provenance}, introduced in G2-prime to close a
loophole exploited by the 7B model in G2).

\subsection{Gate Logic}

The gate emits one of three verdicts:

\begin{description}[leftmargin=*,noitemsep]
  \item[\textsc{admit}] All steps pass provenance, schema-conformance, and
    operator-type checks; the pipeline returns the step chain's output as the
    answer.
  \item[\textsc{reject}] One or more steps fail.  The gate records a failure
    code: \texttt{TYPE\_MISMATCH}, \texttt{PROVENANCE\_SPAN\_NOT\_FOUND},
    \texttt{UNADMITTED\_STEP\_REF}, or \texttt{SCHEMA\_ERROR}.
  \item[\textsc{abstain}] The formaliser produced no parseable schema
    (\texttt{PARSE\_FAIL}) or an empty step list (\texttt{EMPTY\_STEPS}).
\end{description}

For this study, only \textsc{admit} produces an answer.
Both \textsc{reject} and \textsc{abstain} count as no-answer (hence reduce
coverage).

\subsection{Engineering History}

The gate underwent three development rounds before the formal experiment:

\begin{itemize}[leftmargin=*,noitemsep]
  \item \textbf{G1} (abandoned): Initial prototype on a 3B model; formalisation
    too weak to yield reliable step schemas.
  \item \textbf{G2} (killed under §7 SVSS rule): 7B model, text-search
    provenance.  Manual review found 62\% of non-trivial admits had
    correct gate verdicts but wrong semantics (\svss{}), indicating the gate
    was permissive rather than discriminating.
  \item \textbf{G2-prime} (promoted to formal test): 14B model (Qwen2.5-14B),
    v4 few-shot prompt with \emph{bridge\_comparison} and nationality-inference
    examples.  SVSS dropped to $\approx\!15\%$; correct-formalisation rate
    reached $\approx\!50\%$.
\end{itemize}

% ─── 4. Experimental Setup ────────────────────────────────────────────────────
\section{Experimental Setup}
\label{sec:setup}

\subsection{Dataset}

We sample 72 questions from the \emph{decomposable} subset of
2WikiMultihopQA \citep{2wikimultihopqa}, stratified by question type.
The final sample comprises 25 \emph{bridge\_comparison}, 22
\emph{comparison}, and 25 \emph{compositional} questions.

\subsection{Evaluation Protocol}

Both arms run on the same 72 questions.  Answers are manually scored as
\textsc{correct} or \textsc{wrong} by the first author using reference
answers and supporting Wikipedia paragraphs.

\paragraph{Preregistration.}  The evaluation plan was filed in
\texttt{PREREG.md} before any G3 runs.  The promotion criterion is:

\[
\text{full-coverage accuracy}(\text{\armii{}}) -
\text{full-coverage accuracy}(\text{\armi{}}) \geq +15\,\text{pp}
\quad\textbf{AND}\quad p < 0.05
\]

on a one-tailed \mcnem{} test over jointly-answered questions.

\paragraph{Primary metric.}  The \mcnem{} test \citep{mcnemar1947} is applied
to the $n_\text{joint}=32$ questions where \armii{} did not abstain.
Full-coverage accuracy---correct answers divided by total questions---is the
secondary metric reported for each arm.

\paragraph{Coverage and risk.}  Coverage is the fraction of questions for
which an arm produces an answer.  Risk is the error rate among answered
questions: $\text{risk} = \text{wrong} / \text{answered}$.

\subsection{Structurally Valid, Semantically Wrong (\svss{})}

A step schema that passes all gate checks yet answers the wrong question is
flagged as \svss{}.  Under the preregistered §7 kill rule, any pilot with
$>50\%$ \svss{} among non-trivial admits is killed before the formal \mcnem{}
test.  G2 was killed under this rule; G3 passed it (G2-prime SVSS $\approx 15\%$).

% ─── 5. Results ───────────────────────────────────────────────────────────────
\section{Results}
\label{sec:results}

\subsection{Overall Accuracy}

Table~\ref{tab:main} summarises the main results.  \armi{} achieves
80.6\% accuracy at full coverage.  \armii{} admits 32 of 72 questions
(44.4\%) and achieves 84.4\% precision on those admits, but its full-coverage
accuracy is only 37.5\%---a deficit of $-43.1$ pp relative to \armi{}.

\begin{table}[h]
\centering
\caption{Main results on 2WikiMultihopQA decomposable questions ($n=72$).
  \emph{Full-coverage accuracy} counts correct answers over all 72 questions.
  \emph{Precision} counts correct answers over questions admitted (answered).}
\label{tab:main}
\begin{tabular}{lcccc}
\toprule
\textbf{Arm} & \textbf{Coverage} & \textbf{Precision} & \textbf{Full-cov.\ acc.} & \textbf{Abstain} \\
\midrule
\armi{} (CoT)   & 100.0\% (72/72) & 80.6\% & \textbf{80.6\%} & 0 \\
\armii{} (Gate) & 44.4\% (32/72)  & 84.4\% & 37.5\% & 40 (55.6\%) \\
\bottomrule
\end{tabular}
\end{table}

\subsection{\mcnem{} Test}

The \mcnem{} test is applied to the 32 questions admitted by \armii{}.
Table~\ref{tab:mcnemar} shows the 2$\times$2 contingency table.

\begin{table}[h]
\centering
\caption{\mcnem{} contingency table on jointly-answered questions ($n_\text{joint}=32$).
  Discordant cells: $b=4$ (\armi{} only correct), $c=2$ (\armii{} only correct).}
\label{tab:mcnemar}
\begin{tabular}{lcc}
\toprule
 & \textbf{\armii{} Correct} & \textbf{\armii{} Wrong} \\
\midrule
\textbf{\armi{} Correct} & 25 & 4 \\
\textbf{\armi{} Wrong}   & 2  & 1 \\
\bottomrule
\end{tabular}
\vspace{2pt}
\footnotesize{$\chi^2(1)=0.667$, $p=0.41$. 40 abstentions excluded.}
\end{table}

The discordant pair $(b{=}4,\,c{=}2)$ yields $\chi^2(1) = (b-c)^2/(b+c) = 4/6 = 0.667$,
$p = 0.41$.  This is far below the critical value $\chi^2_{0.05}(1) = 3.84$
and fails the preregistered promotion criterion on both grounds.

\subsection{Per-Type Breakdown}

Table~\ref{tab:pertype} shows accuracy broken down by question type.
The gate's behaviour differs sharply across types.

\begin{table}[h]
\centering
\caption{Per-type accuracy. \emph{Cov} = coverage (admitted / total).
  \emph{Prec} = precision when admitted (correct / admitted).
  \emph{Full} = full-coverage accuracy (correct / total).}
\label{tab:pertype}
\begin{tabular}{lcccccc}
\toprule
\textbf{Type} & $n$ & \multicolumn{2}{c}{\armi{}} & \multicolumn{3}{c}{\armii{}} \\
\cmidrule(lr){3-4}\cmidrule(lr){5-7}
              &     & Full   & — & Cov    & Prec   & Full \\
\midrule
bridge\_comparison & 25 & 96.0\% & — & 68.0\% & 88.2\% & 60.0\% \\
comparison         & 22 & 86.4\% & — & 59.1\% & 92.3\% & 54.5\% \\
compositional      & 25 & 60.0\% & — &  8.0\% &  0.0\% &  0.0\% \\
\midrule
\textbf{Overall}   & 72 & 80.6\% & — & 44.4\% & 84.4\% & 37.5\% \\
\bottomrule
\end{tabular}
\end{table}

\emph{bridge\_comparison} and \emph{comparison} show moderately useful
gate precision (88--92\%) but significant coverage loss (32--41\%
abstain).  \emph{compositional} is the dominant failure: 23 of 25 questions
abstain (92\%); the 2 that are admitted are both wrong (0\% precision).

% ─── 6. Analysis ──────────────────────────────────────────────────────────────
\section{Analysis}
\label{sec:analysis}

\subsection{The Abstain Bottleneck}

The gate's primary failure is not incorrect verdicts on the questions it
admits---84.4\% precision is respectable and exceeds \armi{}'s 80.6\%
baseline on the same subset.  The failure is that the gate refuses to admit
most questions.  Of the 40 abstentions, the dominant reason is that the
formaliser either produces no parseable JSON (\texttt{PARSE\_FAIL}) or an
empty step list (\texttt{EMPTY\_STEPS}) for question types it was not
designed for.

This is a known failure mode in neurosymbolic systems: the
formalisation step is itself an LLM generation, subject to all the
fragilities of sequence generation, and the gate is binary (admit or
abstain)---there is no soft fallback.

\subsection{Coverage--Risk Tradeoff}

Figure~\ref{fig:riskcoverage} plots the two arms in coverage--risk space.
\armii{} achieves a marginal risk reduction ($-3.8$ pp: from 19.4\% to
15.6\%) at the cost of a large coverage reduction ($-55.6$ pp: from 100\%
to 44.4\%).  For a full-coverage QA system, this tradeoff is not worthwhile.
An abstaining system that covers only 44\% of queries cannot replace a
general-purpose CoT system.

\begin{figure}[h]
\centering
% \includegraphics[width=0.5\textwidth]{figures/risk_coverage.pdf}
\framebox[0.5\textwidth]{\parbox{0.49\textwidth}{\centering\vspace{20pt}
  [Figure 4: Coverage--Risk scatter]\vspace{20pt}}}
\caption{Coverage--risk tradeoff. Arm~0 (blue circle): coverage=100\%,
  risk=19.4\%. Arm~2 (amber circle): coverage=44.4\%, risk=15.6\%.
  The gate gains $-3.8$ pp risk at the cost of $-55.6$ pp coverage.}
\label{fig:riskcoverage}
\end{figure}

\subsection{Compositional Type: Total Gate Failure}

The \emph{compositional} type (questions requiring multi-step chains not
directly supported by a single operator) completely defeats the gate.
The formaliser was trained primarily on \emph{bridge\_comparison}
examples---a natural-language inference pattern where two entity
attributes are compared.  Compositional questions require chaining
\texttt{conjoin} with \texttt{arithmetic} or \texttt{temporal\_order}
in ways not covered by the few-shot examples, leading to nearly universal
parse failure.

This confirms a key limitation of few-shot formalisation: coverage is
only as broad as the example distribution.

\subsection{Mode A vs.\ Mode B}

An earlier internal study (\emph{Mode~A}, formalization-only, no QA
grounding) used a 92\% bridge\_comparison sample, inflating gate
performance estimates.  The formal G3 experiment (\emph{Mode~B}) samples
the full decomposable distribution and finds only 68\% bridge\_comparison.
The discrepancy accounts for most of the observed coverage decline from
internal estimates to the formal result, and underscores the importance of
representative sampling in pilot studies.

% ─── 7. Discussion ────────────────────────────────────────────────────────────
\section{Discussion}
\label{sec:discussion}

\subsection{Why Hard Gates Fail on Generated Formalisms}

The fundamental tension is between two requirements that pull in opposite
directions.  A hard gate needs a \emph{reliable formaliser}: if the
formaliser often fails to produce a valid schema, most questions abstain and
coverage collapses.  But a reliable formaliser on natural-language multi-hop
questions---which may involve arbitrary cultural, temporal, geographic, or
numeric reasoning---requires either a very large model (reducing the
efficiency argument) or a very constrained question type (reducing generality).

Neurosymbolic successes (Logic-LM, SatLM, LINC) mostly operate in domains
with well-defined, well-modelled logical structure: arithmetic word problems,
formal proofs, boardgame rules.  Multi-hop QA from Wikipedia spans informal,
compositional, culturally-specific reasoning for which no clean symbolic
vocabulary exists.

\subsection{The Role of Error Correlation}

Both arms derive from the same Qwen2.5-14B model family.  As
\citet{correlatedErrors} and \citet{llmsJury} document, models in the same
family make correlated errors; the \mcnem{} concordance table (25 both-right,
1 both-wrong, $n{=}32$) is consistent with this.  A gate built on a diverse,
independent verifier might see fewer concordant-correct cells and more
discordant gains, potentially improving the \mcnem{} outcome.  However,
the dominant failure---abstention, not wrong admission---is structural
and cannot be addressed by reducing error correlation alone.

\subsection{When Could a Gate Help?}

Based on our results, a typed gate is likely to help when:
\begin{enumerate}[leftmargin=*,noitemsep]
  \item The question type is \emph{predictable} and covered by the formaliser's
    training distribution (our bridge\_comparison and comparison results hint at
    this: 88--92\% precision when admitted).
  \item An \emph{abstain} outcome is acceptable (selective QA, human-in-the-loop).
  \item The formaliser is a much larger or more capable model than the answerer.
\end{enumerate}

For general-coverage QA systems, the evidence suggests that a soft verifier
(a learned scoring model) or a cascade approach
(gate $\rightarrow$ fallback to CoT on abstain) would better preserve coverage.

% ─── 8. Limitations ───────────────────────────────────────────────────────────
\section{Limitations}
\label{sec:limitations}

\paragraph{Sample size.}  $n{=}72$ is modest; the \mcnem{} test has limited
power, and the per-type cells ($n \in \{22, 25\}$) are too small for
subgroup inference.

\paragraph{Single model family.}  Both arms use Qwen2.5-14B or related
variants.  Results may differ with a larger or more diverse formaliser.

\paragraph{Manual scoring.}  Scores were assigned by a single annotator.
Inter-annotator agreement was not measured.

\paragraph{Preregistration scope.}  The preregistration was internal
(\texttt{PREREG.md} in the project repository), not filed with an external
registry.  The \mcnem{} primary metric and $+15$ pp threshold were locked
before runs, but the §7 SVSS kill rule was added after G2 based on
observed failure modes.

\paragraph{Operator coverage.}  The four operators were chosen for the
decomposable subset; other question types (bridge-only, single-hop) are
not evaluated.  The 92\% abstain rate on compositional questions likely
reflects a systematic gap rather than random formaliser error.

\paragraph{No cascade baseline.}  We do not evaluate a gate+fallback
(CoT when gate abstains) design.  The paper measures the gate in isolation;
a cascade might recover coverage at some precision cost.

% ─── 9. Conclusion ────────────────────────────────────────────────────────────
\section{Conclusion}
\label{sec:conclusion}

We find that a typed-operator gate does not improve on direct
chain-of-thought generation on 2WikiMultihopQA decomposable questions.
The gate's precision is competitive (84.4\%) on the questions it admits,
but it abstains on 56\% of questions, collapsing full-coverage accuracy
from 80.6\% to 37.5\%.  The \mcnem{} test is non-significant
($\chi^2(1) = 0.667$, $p = 0.41$), failing both required promotion criteria.

The failure is predominantly an abstain bottleneck, driven by the gap
between the formaliser's training distribution (bridge\_comparison-heavy)
and the full decomposable question distribution (including compositional
types).  For general-coverage multi-hop QA, the evidence favours softer
verification or cascade designs over hard binary gates.

The \latch{} pilot provides a negative-result baseline for future
neurosymbolic multi-hop QA systems on 2WikiMultihopQA, and demonstrates
that a preregistered pilot protocol can efficiently kill unpromising
designs without large-scale compute.

% ─── Acknowledgements ─────────────────────────────────────────────────────────
\section*{Acknowledgements}

The 2WikiMultihopQA dataset is released under the Apache~2.0 license by
\citet{2wikimultihopqa}.  Experiments were run locally on consumer hardware
(NVIDIA RTX 3090) using Ollama and the publicly available Qwen2.5-14B
weights.

% ─── References ───────────────────────────────────────────────────────────────
\bibliographystyle{plainnat}
\bibliography{paper}

\end{document}
